\begin{textsample}{NEG Dim 2 – 2022\_c5 – Score -1.00 – 2022\_c5/enade\_2022\_tec\_gest\_qual\_854a44b817.txt}  \label{ex:f2_neg_219}
CAPÍTULO 6 DISTRIBUIÇÃO DE CONCEITOS A finalidade deste capítulo é apresentar um panorama nacional da distribuição do Conceito Enade 2022.
Para traçar esse panorama, recorre -se a tabelas e gráficos, que são acompanhados de comentários e análises que contribuem para a contextualização dos dados neles expostos.
Para visualizar o panorama do Enade/2002, em nível nacional, na primeira seção do capítulo, os dados sobre o número de cursos participantes são organizados por grande região.
Na segunda seção, o mapeamento do conceito do exame é feito por categoria administrativa das IES e modalidade de oferta dos cursos, de maneira cruzada às grandes regiões.
Na última seção, a lógica de apresentação dos dados é feita por organização acadêmica e, também, grande região de funcionamento dos cursos, no país como um todo. 6.1 PANORAMA NACIONAL DA DISTRIBUIÇÃO DOS CONCEITOS A lógica adotada na apresentação dos dados, nesta seção e nas duas seguintes, pressupõe o registro de duas informações de natureza reiterativa.
A primeira informação diz respeito aos valores possíveis do Conceito Enade, que variam em níveis de 1 a 5, sendo que, à medida que esses valores aumentam, o curso recebe uma avaliação melhor, além disso, os cursos com somente um ou sem nenhum concluinte participante não obtêm o Conceito Enade, ficando Sem Conceito ( SC ), como já ressaltado neste relatório.
O segundo registro é que eventual diferença entre os cursos tabulados neste capítulo e no Capítulo 3 corresponde a cursos que não foram avaliados, ou seja, aqueles, em princípio, sem estudantes concluintes que tivessem sido inscritos para a prova.
Na Tabela 6.1, são apresentadas a quantidade e a distribuição de cursos de Tecnologia em Gestão da Qualidade participantes do Enade/2022, por faixa de conceito e grande região.
Observando -se os dados da Tabela 6.1, nota -se que, dos 112 cursos participantes, 38 ( 33,9 \% ) classificaram -se com conceito 3, o valor modal.
Esse foi também o conceito modal na maioria das regiões : Nordeste ( 71,4 \% ), Sudeste ( 30 \% ) e Sul ( 29,6 \% ).
O conceito 2 foi o segundo mais frequente em nível nacional ( 17,0 \%, correspondendo a 19 cursos ).
O conceito 4 foi o terceiro, com 15 cursos ( 13,4 \% ).
Houve, ainda, 4 cursos ( 3,6 \% ) que receberam o conceito 5, e 2 cursos ( 1,8 \% ) que receberam o conceito 1.
Dos 112 cursos de Tecnologia em Gestão da Qualidade, 34 ( 30,4 \% ) ficaram Sem Conceito ( SC ).
A região Norte participou com 8 cursos.
Desses, 3 cursos receberam o conceito 2, o que equivale a 37,5 \% do total regional, o conceito modal nessa região.
Os conceitos 1, 4 ou 5 não foram atribuídos a nenhum curso e o conceito 3 foi atribuído a 2 cursos ( 25 \% ).
Ainda nessa região, 3 cursos ( 37,5 \% ) ficaram Sem Conceito ( SC ).
A região Nordeste participou com 14 cursos.
Nessa região, 10 cursos ( 71,4 \% em termos regionais ) obtiveram conceito 3, o conceito modal para a região.
Os conceitos 1 e 2 foram atribuídos, respectivamente, a 1 curso ( 7,1 \% ) e a 2 cursos ( 14,3 \% ).
Ainda nessa região, 1 curso ( 7,1 \% ) foi avaliado com o conceito 4 e nenhum curso com o conceito 5 ou Sem Conceito ( SC ).
Para os 60 cursos participantes da região Sudeste, o conceito 3 foi o conceito modal, como já anteriormente assinalado, tendo sido obtido por 18 cursos ( 30,0 \% ).
O conceito 1 foi atribuído a um curso ( 1,7 \% ) ; o conceito 2, a 7 cursos ( 11,7 \% ) ; e os conceitos 4 e 5, a 8 cursos ( 13,3 \% ) e a 2 cursos ( 3,3 \% ), respectivamente.
Nessa região, 24 cursos ficaram Sem Conceito ( SC ).
Em relação aos 27 cursos da região Sul, houve predominância do conceito 3 ( conceito modal ), atribuído a 29,6 \% dos cursos da região, ou seja, atribuído a 8 dos 27 cursos participantes na região Sul.
Nenhum curso recebeu o conceito 1 ; o conceito 2 foi atribuído a 6 cursos ( 22,2 \% ) ; e os conceitos 4 e 5, a 6 ( 22,2 \% ) e a 2 cursos ( 7,4 \% ), respectivamente.
Nessa região, cinco cursos ficaram Sem Conceito ( SC ).
Já dos 3 cursos participantes na região Centro-Oeste, 1 ( 33,3 \% em termos regionais ) recebeu conceito 2 e 2 cursos ficaram Sem Conceito ( SC ).
No gráfico de distribuição cumulativa de cada grande região mostrado a seguir, para cada valor da abscissa x ( que corresponde aos valores possíveis para o Conceito Enade ), o valor apresentado no eixo y é a porcentagem de cursos da grande região que obtiveram Conceito Enade menor ou igual a x.
A partir dessas diretrizes, no Gráfico 6.1, é apresentada a distribuição cumulativa do Conceito Enade, segundo grande região.
Nesse gráfico, poligonais mais à esquerda representam valores mais baixos do Conceito Enade, e poligonais mais à direita, valores mais altos. 6.2 CONCEITOS POR CATEGORIA ADMINISTRATIVA, MODALIDADE DE OFERTA DOS CURSOS E GRANDE REGIÃO Na Tabela 6.2, é apresentada a distribuição dos cursos de Tecnologia em Gestão da Qualidade participantes do Enade/2022, por categoria administrativa e modalidade de oferta, de acordo com os conceitos por eles alcançados, segundo as grandes regiões brasileiras.
De acordo com as informações da Tabela 6.2, dos 112 cursos participantes, 7 eram ministrados em instituições públicas, e 105, em instituições privadas.
Quanto à modalidade de oferta, 63 cursos eram a distância, e 49 eram presenciais.
Ainda conforme as informações da mesma tabela, em termos nacionais, entre instituições públicas e privadas, observa -se que, 4 cursos foram avaliados com conceito 5.
Dos 7 cursos participantes de IES públicas, os conceitos 3, 4 e 5 foram atribuídos a 2 cursos cada.
Além disso, 1 curso de instituição pública participante foi avaliado com o conceito 2.
Nessa categoria adminsitrativa, nenhum curso ficou Sem Conceito ( SC ).
Na rede privada, constituíram o conceito modal foi o 3, atribuído a 36 cursos.
Dos demais cursos participantes, 2 receberam o conceito 1 ; 18 receberam o conceito 2 ; 13 receberam o conceito 4 ; e 2 receberam o conceito 5.
Nessa categoria administrativa, 34 cursos ficaram Sem Conceito ( SC ).
Dos 49 cursos participantes da modalidade presencial, o conceito 3 foi o valor modal, atribuído a 11 cursos.
Dos demais cursos nessa modalidade de oferta, 2 cursos receberam o conceito 1 ; 6 cursos receberam o conceito 2 ; 7 cursos, o conceito 4 ; e 3 cursos, o conceito 5.
Nessa modalidade de oferta, 20 cursos ficaram Sem Conceito ( SC ).
Dos 63 cursos participantes da modalidade a distância, o conceito 3 foi o valor modal, atribuído a 27 cursos.
Dos demais cursos nessa modalidade de oferta, nenhum recebeu o conceito 1 ; 13 cursos receberam o conceito 2 ; 8 cursos, o conceito 4 ; e 1 curso, o conceito 5.
Nessa modalidade de oferta, 14 cursos ficaram Sem Conceito ( SC ).
Na análise por região, observa -se que as instituições privadas concentraram o total de cursos participantes ( 8 ) da região Norte.
Desses, 2 cursos obtiveram conceito 3 e 3 cursos, o conceito 2 ( conceito modal ).
Nessa combinação de categoria administrativa e grande região, nenhum curso recebeu os conceitos 1, 4 ou 5 ; e 3 cursos ficaram Sem Conceito ( SC ).
Na região Nordeste, a rede privada concentrou 11 dos 14 cursos participantes.
Entre os cursos oferecidos por IES privadas nessa região, 9 cursos foram avaliados com conceito 3.
Nenhum curso obteve os conceitos 4 ou 5 ; 1 curso obteve o conceito 1 ; e 1, o conceito 2.
Nessa categoria administrativa e grande região, nenhum curso ficou Sem Conceito ( SC ).
As instituições públicas da região Nordeste participaram com 3 cursos, que obtiveram os conceitos 2, 3 e 4.
Nenhum curso ficou Sem Conceito ( SC ) nessa região e categoria administrativa.
Dos cursos oferecidos na região Nordeste, 4 eram presenciais ; nenhum deles ficou Sem Conceito ( SC ) ; 1 curso recebeu o conceito 1 ; 1 curso, o conceito 2 ; 1 curso, o conceito 3 ; e 1, o conceito 4.
Na modalidade a distância, participaram 10 cursos, tendo 1 deles obtido o conceito 2 ; e 9, o conceito 3.
Nessa modalidade e região, nenhum curso ficou Sem Conceito ( SC ) ou recebeu conceitos 1, 4 ou 5.
Na região Sudeste, a proporção de cursos da rede privada correspondeu a 58 dos 60 cursos participantes.
Nessa combinação de categoria administrativa e grande região, o conceito modal foi o 3, atribuído a 17 cursos.
Observa -se que 24 cursos ficaram Sem Conceito ( SC ).
Os demais foram avaliados com os seguintes conceitos : 1 ( 1 curso ), 2 ( 7 cursos ), 4 ( 7 cursos ) e 5 ( 2 cursos ).
Dos 2 cursos oferecidos por instituições públicas na região Sudeste, 1 recebeu o conceito 3 e 1, o conceito 4.
Nessa combinação de categoria administrativa e grande região, nenhum curso ficou Sem Conceito ( SC ), e nenhum recebeu conceito 1, 2 ou 5.
Também na região Sudeste, 30 cursos foram oferecidos na modalidade presencial e 30 na modalidade a distância.
Para os cursos presenciais, o conceito 1 foi atribuído a 1 curso ; o conceito 2 foi atribuído a 2 cursos ; o conceito 3 foi atribuído a 8 cursos ; o conceito 4, a 2 cursos ; e o 5 a 2 cursos.
Ficaram Sem Conceito ( SC ) 16 cursos presenciais na região Sudeste.
Além disso, nessa região, 30 cursos também foram oferecidos na modalidade a distância.
Desses, 1 curso recebeu o conceito 5 ; 5 cursos receberam conceito 2 ; 6 cursos, o conceito 4 ; 10 cursos, o conceito 3, o valor modal ; e 8 cursos ficaram Sem Conceito ( SC ).
As instituições privadas concentraram 25 dos 27 cursos participantes da região Sul.
Desses, 8 cursos obtiveram conceito 3, o conceito modal.
Nessa combinação de categoria administrativa e grande região, nenhum curso recebeu o conceito 1 ; 6 receberam o conceito 2 ; 6, o conceito 4 ; nenhum curso recebeu o conceito 5 ; e 5 cursos ficaram Sem Conceito ( SC ).
As instituições públicas na região Sul participaram com apenas 2 cursos, aos quais foram atribuídos o conceito 5.
Na região Sul, 12 cursos foram oferecidos na modalidade presencial.
Desses, 4 cursos ficaram Sem Conceito ( SC ) ; nenhum recebeu conceito 1 ; 1 curso recebeu conceito 2 ; 1 curso recebeu conceito 3 ; 4 cursos receberam conceito 4 ; e 2, conceito 5.
Na modalidade a distância, foram oferecidos 15 cursos.
Desses, 1 curso ficou Sem Conceito ( SC ) ; nenhum recebeu conceito 1 ; 5 cursos receberam conceito 2 ; 7 cursos receberam conceito 3 ; 4 cursos receberam conceito 4 ; e nenhum curso recebeu conceito 5.
Na região Centro-Oeste, os 3 cursos participantes eram de IES privadas.
Desses, 1 curso recebeu conceito 2 e 2 cursos ficaram Sem Conceito ( SC ).
Nessa região, todos os cursos ofertados são da modalidade a distância.
A título de síntese, no Gráfico 6.2, apresenta -se a distribuição cumulativa do Conceito Enade, segundo a categoria administrativa da IES.
Nesse gráfico, como dito anteriormente, poligonais mais à esquerda representam valores mais baixos do Conceito Enade, e poligonais mais à direita, valores mais altos.
É importante observar que, para cada poligonal presente nesse gráfico, a porcentagem constante no eixo y se refere ao universo de cursos oferecidos por instituições públicas ( poligonal vermelha ) ou privadas ( poligonal azul ). 6.3 CONCEITOS POR ORGANIZAÇÃO ACADÊMICA E GRANDE REGIÃO Na Tabela 6.3, encontra -se a distribuição dos conceitos atribuídos aos cursos participantes do Enade/2022, na área de Tecnologia em Gestão da Qualidade, por organização acadêmica, segundo as grandes regiões brasileiras.
Como mostram os dados da Tabela 6.3, dos 112 cursos de Tecnologia em Gestão da Qualidade participantes, 41 eram oferecidos em universidades, 54 em centros universitários, 14 em faculdades e 3 em CEFET/IF.
De acordo com os dados apresentados na mesma tabela, de todos os 4 cursos avaliados com o conceito 5, 2 eram vinculados a universidades.
Os cursos vinculados a esse tipo de organização acadêmica tiveram o conceito 3 como conceito modal, atribuído a 17 cursos.
Os demais cursos ofertados por universidades receberam os conceitos 2 ( 9 cursos ) e 4 ( 4 cursos ).
Nenhum curso recebeu conceito 1 e 9 cursos ficaram Sem Conceito ( SC ).
Dos cursos vinculados a centros universitários, o conceito modal também foi o 3, atribuído a 16 cursos.
Observa -se que 21 cursos vinculados a esse tipo de organização acadêmica ficaram Sem Conceito ( SC ), e os outros cursos receberam os conceitos 2 ( 9 cursos ), 4 ( 7 cursos ) e 5 ( 1 curso ).
Dos 14 cursos mantidos por faculdades, 4 receberam conceito 3, o conceito modal.
O segundo conceito mais frequentemente atribuído aos cursos foi o conceito 4, recebido por 3 cursos.
Na sequência, 2 cursos obtiveram conceito 1 e 1 curso, o conceito 2.
Nenhum curso recebeu o conceito 5 e 4 cursos ficaram Sem Conceito ( SC ).
Dos 3 cursos mantidos por CEFET/IF, um recebeu conceito 3 ; um recebeu conceito 4 e um recebeu o conceito 5.
Nenhum curso ficou Sem Conceito ( SC ).
Considerando -se, separadamente, as regiões brasileiras, verifica -se que, na região Norte, as universidades participaram com 4 cursos, dos quais 2 receberam o conceito 3 e 2 cursos ficaram Sem Conceito ( SC ).
Os centros universitários da região Norte foram representados por 3 cursos, dos quais 2 obtiveram conceito 2, o valor modal.
Os conceitos 1, 3, 4 e 5 não foram atribuídos a nenhum curso.
Um curso ficou Sem Conceito ( SC ).
As faculdades participaram com 1 curso na região Norte, o qual recebeu o conceito 2.
Nenhum curso ficou Sem Conceito ( SC ).
Além disso, nessa região e organização acadêmica, nenhum curso foi oferecido em CEFET/IF.
Na região Nordeste, as universidades participaram com 6 dos 14 cursos da área de Tecnologia em Gestão da Qualidade oferecidos.
O conceito modal foi 3, atribuído a 10 cursos.
Um curso obteve o conceito 4 ; 1 curso, o conceito 1 ; e outros 2, o conceito 2.
Nenhum curso oferecido por universidades no Nordeste recebeu o conceito 5 ou ficou Sem Conceito ( SC ).
Os centros universitários contaram com 6 cursos participantes na região Nordeste e todos receberam o conceito 3.
As faculdades foram representadas por 1 curso na região Nordeste, o qual recebeu conceito 1.
Os CEFET/IF foram representados por 1 curso que recebeu o conceito 4.
Na região Sudeste, as universidades concentraram 21 dos 60 cursos de Tecnologia em Gestão da Qualidade da região.
Dos cursos oferecidos em IES com este tipo de organização acadêmica, o conceito modal foi 3, atribuído a 8 cursos.
Os demais cursos receberam os conceitos 2 ( 5 cursos ) ; 4 ( 1 curso ) ; 5 ( 1 curso ) ; e 6 cursos ficaram Sem Conceito ( SC ).
Os centros universitários participaram com 29 cursos na região Sudeste, dos quais 6 obtiveram o conceito 4.
O conceito 2 foi atribuído a dois cursos ; o conceito 3, a 5 cursos ; o conceito 5, a 1 curso.
Nessa combinação de organização acadêmica e região, 15 cursos ficaram Sem Conceito ( SC ).
Já as faculdades foram representadas por 9 cursos na região Sudeste, dos quais três ficaram Sem Conceito ( SC ).
O conceito 3 foi atribuído a 4 cursos.
Aos demais cursos foram atribuídos os conceitos 1 ( 1 curso ) e 4 ( 1 curso ).
O único curso participante de CEFET/IF, na região Sudeste, o conceito foi 3 e nenhum curso ficou Sem Conceito ( SC ).
Dos 27 cursos da região Sul, 10 eram mantidos por universidades.
Desses, 3 cursos obtiveram o conceito 3 ; 3 cursos, o conceito 4 ; 2 cursos, o conceito 2 ; 1 curso, o conceito 5 e nenhum curso recebeu o conceito 1.
Somente 1 curso ficou Sem Conceito ( SC ) nessa combinação de organização acadêmica e região.
Dos 13 cursos participantes de centros universitários da região Sul, três ficaram Sem Conceito ( SC ), e 5 cursos receberam conceito 3, o valor modal.
Os conceitos restantes ficaram assim \textbf{distribuídos} : 4 cursos ( conceito 2 ) ; 1 curso ( conceito 4 ).
Quanto aos 3 cursos vinculados a faculdades na região Sul, 2 cursos receberam o conceito 4 ( valor modal ) e somente um curso ficou Sem Conceito ( SC ).
O único curso oferecido em CEFET/IF foi avaliado com o conceito 5.
Nesse tipo de organização, nenhum curso da região Sul ficou Sem Conceito ( SC ) ou com os conceitos 1, 2, 3 e 4.
Na região Centro-Oeste, os 3 cursos oferecidos eram mantidos por centros universitários.
Para esse tipo de organização acadêmica, 1 curso recebeu o conceito 2 e 2 cursos ficaram Sem Conceito ( SC ).
De maneira semelhante à forma como se procedeu, em relação à seção anterior, encerra -se esta seção com o Gráfico 6.3, que apresenta a distribuição do Conceito Enade dos cursos de Tecnologia em Gestão da Qualidade segundo a organização acadêmica da IES.
A leitura desse último gráfico é semelhante à leitura dos anteriores, seguindo -se a lógica de que poligonais deslocadas para a esquerda correspondem a valores mais baixos do Conceito Enade ( 1, 2 e 3 ) e poligonais deslocadas para a direita correspondem a valores mais altos do Conceito Enade.

% matched lemmas: distribuir
\end{textsample}
